\chapter*{Notation und Begriffsdefinitionen}
\addcontentsline{toc}{chapter}{Notation und Begriffsdefinitionen}

In diesem Dokument werden die folgenden typografischen Konventionen verwendet:

\begin{description}
    \itemsep0.5em
    \item[\texttt{Nichtproportionale Schrift}] bezeichnet konkrete Klassen, Strukturen, Methoden und Code-Artefakte in \textit{helios} (z.\,B. \texttt{GameLoop}, \texttt{SparseSet}, \texttt{UpdateContext}).

    \item[\textit{Kursive Schrift}] kennzeichnet allgemeine Fachbegriffe, Konzepte und englischsprachige Termini (z.\,B. \textit{Composition over Inheritance}, \textit{Data-Oriented Design}, \textit{Spawning}).
\end{description}

Da manche Bezeichnungen zugleich einen allgemeinen Fachbegriff und eine konkrete Klasse benennen, wie etwa \textit{GameLoop} als architektonisches Konzept und \texttt{GameLoop} als Klasse in \textit{helios}, ergibt sich die jeweilige Lesart aus dem Kontext.\\

Der Begriff \textit{System} wird in diesem Dokument sowohl für Laufzeit- und Subsysteme der Engine (z.\,B. Rendering- oder Kollisionssysteme) als auch für \textit{ECS-Systeme} im Sinne einer \textit{Entity-Component-System}-Architektur verwendet.
Die jeweils gemeinte Bedeutung ergibt sich aus dem Kontext.\\

Der Begriff \textit{Array} wird in diesem Dokument stellvertretend auch für \texttt{std::vector} verwendet, sofern eine Unterscheidung für den jeweiligen Kontext nicht relevant ist.\\

An einigen Stellen wird bei generischen Typen auf die explizite Angabe von Template-Parametern verzichtet, sofern diese für das Verständnis des jeweiligen Kontexts nicht erforderlich ist.\\

Ergänzende Informationen werden im Dokument in farblich abgesetzten Textboxen dargestellt:

\vspace{5mm}
\begin{tcolorbox}[colback=black!5, colframe=black, title={Beispiel}]
    Inhalte in dieser Form vertiefen den umgebenden Abschnitt und liefern weiterführende Informationen.
\end{tcolorbox}
\vspace{5mm}

In diesem Dokument wird wiederholt auf das \textit{Prototyp-Dokument}~\cite{SH25} Bezug genommen, das die erste Projektphase und den Entwurf des Prototypen beschreibt.
Das vorliegende Dokument wird im Folgenden als \textit{Engine-Dokument} bezeichnet und dokumentiert die Weiterentwicklung vom Prototypen zur ECS-basierten Engine.\\


Weitere Begriffserklärungen finden sich im Glossar in Kapitel~\ref{ch:glossar}.